% Original tug-of-war design, with rows populated by reproduce.py.
% Arrow geometry uses unrounded results; displayed values use manuscript rounding.
\begin{figure}[t]
\centering
\vspace{1mm}
\usetikzlibrary{arrows.meta}
\definecolor{tugblue}{HTML}{2B6CB0}
\definecolor{tugred}{HTML}{C0453C}
\definecolor{tugink}{HTML}{253746}
\definecolor{tugrope}{HTML}{BCAA8F}
\resizebox{\columnwidth}{!}{%
\begin{tikzpicture}[x=1cm,y=1cm,font=\footnotesize]
  % Light grouping panels keep the original rope, arrows, and numerical geometry.
  \fill[black!2,rounded corners=3pt] (-.12,-.25) rectangle (8.68,-2.38);
  \fill[black!2,rounded corners=3pt] (-.12,-2.40) rectangle (8.68,-4.53);
  \fill[black!2,rounded corners=3pt] (-.12,-4.55) rectangle (8.68,-6.68);
  \def\xL{1.50}
  \def\xR{7.30}
  \def\xC{4.40}
  \def\sK{7.3}
  \def\stealth{-{Stealth[length=2.0mm,width=1.6mm]}}

  % #1 y; #2 protocol; #3 compression; #4 net loss; #5--7 display labels.
  \newcommand{\tugrow}[7]{%
    \node[anchor=east] at (1.42,#1) {#2};
    \draw[tugrope,line width=1.6pt,line cap=round] (\xL,#1) -- (\xR,#1);
    \filldraw[fill=tugblue,draw=none] (\xL,#1) circle (0.055);
    \filldraw[fill=tugred,draw=none] (\xR,#1) circle (0.055);
    \pgfmathsetmacro{\lenb}{abs(#3)*\sK}
    \ifdim#3pt>0pt
      \draw[tugblue,\stealth,line width=1.7pt,line cap=round]
        (\xL+\lenb,#1) -- (\xL,#1);
    \else
      \draw[tugblue,\stealth,line width=1.7pt,line cap=round]
        (\xL,#1) -- (\xL+\lenb,#1);
    \fi
    \pgfmathsetmacro{\lenr}{abs(#4)*\sK}
    \ifdim#4pt>0pt
      \draw[tugred,\stealth,line width=1.7pt,line cap=round]
        (\xR-\lenr,#1) -- (\xR,#1);
    \else
      \draw[tugred,\stealth,line width=1.7pt,line cap=round]
        (\xR,#1) -- (\xR-\lenr,#1);
    \fi
    \pgfmathsetmacro{\kx}{\xC+(#4-(#3))*\sK}
    \filldraw[fill=tugink,draw=white,line width=0.7pt] (\kx,#1) circle (0.085);
    \node[anchor=west,text=tugblue] at (\xL+0.04,#1+0.30) {#5};
    \node[anchor=east,text=tugred] at (\xR-0.04,#1+0.30) {#6};
    \node[anchor=east] at (8.55,#1+0.10) {#7};
  }%
  \newcommand{\tughead}[3]{%
    \node[anchor=west,font=\footnotesize\bfseries] at (0,#1) {#2};
    \node[anchor=east,font=\footnotesize\itshape,text=black!55] at (8.55,#1) {#3};
  }%
  \newcommand{\tugzero}[2]{%
    \draw[dashed,black!35,line width=0.5pt] (\xC,#1) -- (\xC,#2);
  }%

  \node[anchor=west,text=tugblue] at (\xL,0) {compression};
  \node[anchor=east,text=tugred] at (\xR,0) {net loss};
  \node[anchor=east] at (8.55,0) {net $\Delta A$};
  \node[anchor=center,text=black!60] at (\xC,0) {centre: net $=0$};

  \tughead{-0.45}{MATH-500}{the two pulls nearly cancel}
  \tugzero{-0.70}{-2.30}
@@MATH_ROWS@@

  \tughead{-2.60}{MMLU-Pro}{selection offsets coverage loss}
  \tugzero{-2.85}{-4.45}
@@MMLU_ROWS@@

  \tughead{-4.75}{GSM8K}{both pulls reduce accuracy}
  \tugzero{-5.00}{-6.60}
@@GSM_ROWS@@

  \begin{scope}[shift={(0,-7.02)}]
    \draw[tugblue,\stealth,line width=1.7pt] (1.12,0) -- (0.52,0);
    \node[anchor=west,text=tugblue] at (1.22,0) {compression};
    \draw[tugred,\stealth,line width=1.7pt] (3.02,0) -- (3.62,0);
    \node[anchor=west,text=tugred] at (3.72,0) {net loss};
    \filldraw[fill=tugink,draw=white,line width=0.7pt] (5.52,0) circle (0.085);
    \node[anchor=west] at (5.68,0) {net $\Delta A$};
  \end{scope}
  \node[anchor=west,text=black!60] at (0,-7.44)
    {Arrow length = term size; knot left of centre = net gain.};
\end{tikzpicture}%
}
\caption{\Diagnose{} discussion as a tug-of-war: regret compression versus net ceiling loss (\%).}
\label{fig:tradeoff}
\end{figure}
